\documentclass[11pt,a4paper]{article}
\usepackage{shared/zoocover}
\usepackage[utf8]{inputenc}
\usepackage[T1]{fontenc}
\usepackage{amsmath,amssymb,amsthm}
\usepackage{graphicx}
\usepackage{hyperref}
\usepackage{cleveref}
\usepackage{algorithm}
\usepackage{algorithmic}
\usepackage{booktabs}
\usepackage{multirow}

\newtheorem{theorem}{Theorem}
\newtheorem{lemma}{Lemma}
\newtheorem{definition}{Definition}
\newtheorem{proposition}{Proposition}
\newtheorem{corollary}{Corollary}

\title{\textbf{Z-Chain: A Universal Zero-Knowledge Proof Platform\\with Post-Quantum Security on Zoo Network}}

\author{
Antje Worring, Zach Kelling\thanks{research@zoo.ngo}\\
\textit{Zoo Labs Foundation (501c3)}
}

\date{February 2026}

\begin{document}
\zoocoverpage
\maketitle

\begin{abstract}
We present Z-Chain, the zero-knowledge proof substrate of Zoo Network, providing universal proving, cheap verification, and composable proofs across all Zoo chains. Z-Chain implements a dual-lane architecture: a \emph{production lane} (Groth16, PLONK, STARK) for stable deployments, and a \emph{research lane} for experimental proof systems with versioned capability identifiers. The core primitive is the \emph{Universal Receipt}---a chain-native proof-of-proof object that enables recursive composition, cross-chain verification, and long-term auditability. We integrate post-quantum security through Poseidon2 hash commitments (replacing Pedersen, achieving 2000$\times$ speedup at 7.5$\times$ lower gas cost), ML-DSA signature verification within ZK circuits, and STARK proofs over the Goldilocks field for transparent, quantum-resistant proving. The PQZ (Post-Quantum Zero-Knowledge) algorithm combines these primitives into a unified framework where: (1) proofs are generated using STARKs with Poseidon2 commitments, (2) verification uses ML-DSA-signed receipts, (3) cross-chain export uses Groth16 wrappers for EVM compatibility, and (4) private computation uses FHE-ZK hybrid protocols on Zoo's T-Chain. We benchmark Z-Chain at 12,000 proof verifications per second on GPU-accelerated validators with 1.8-second end-to-end latency from proof submission to receipt finalization. Formal verification in Lean 4 proves soundness, completeness, and zero-knowledge properties for all three production proof systems. This work establishes the first post-quantum zero-knowledge infrastructure integrated natively into a blockchain consensus layer.
\end{abstract}

\section{Introduction}

Zero-knowledge proofs enable a prover to convince a verifier of a statement's truth without revealing the underlying witness. In blockchain systems, ZKPs serve three critical roles: scaling (rollup compression), privacy (confidential transactions), and verifiable computation (off-chain execution with on-chain verification). However, existing ZK infrastructure suffers from three limitations:

\begin{enumerate}
\item \textbf{Quantum vulnerability.} Groth16 and PLONK rely on elliptic curve pairings (BN254, BLS12-381) that are broken by Shor's algorithm on a sufficiently large quantum computer. As NIST has standardized post-quantum algorithms (ML-KEM, ML-DSA, SLH-DSA), ZK systems must migrate to quantum-resistant primitives.

\item \textbf{Fragmentation.} Each proof system (Groth16, PLONK, STARK, Nova, Halo2) produces incompatible proof formats. Cross-system composition requires manual bridging. No standard receipt format exists for proof-of-proof attestation.

\item \textbf{Centralized infrastructure.} Proof generation is computationally intensive, concentrating proving power in well-funded operators. Decentralized proving markets exist (=nil;, Gevulot) but lack native blockchain integration.
\end{enumerate}

Z-Chain addresses all three through a chain-native ZK substrate that is post-quantum by default, universally composable via receipts, and decentralized through Zoo Network's Proof of AI validator set.

\section{Architecture}

\subsection{Chain Placement}

Z-Chain operates as a specialized chain within the Zoo Network multi-chain architecture:

\begin{itemize}
\item \textbf{C-Chain} (Contract Chain): Application deployment. Consumes Z-Chain proofs via precompiles.
\item \textbf{Z-Chain} (ZK Substrate): Proof registration, verification, and receipt generation. This paper's focus.
\item \textbf{T-Chain} (Threshold Chain): FHE decryption coordination for private ZK-FHE hybrid protocols.
\item \textbf{K-Chain} (Key Chain): Post-quantum key management (ML-KEM, ML-DSA) for proof signing.
\end{itemize}

\subsection{Two-Lane Design}

\subsubsection{Production Lane}

Three proof systems are permanently supported with stable verifier contracts:

\begin{table}[h]
\centering
\begin{tabular}{llll}
\toprule
\textbf{System} & \textbf{Setup} & \textbf{PQ-Safe} & \textbf{Verify Gas} \\
\midrule
Groth16 & Trusted & No\textsuperscript{$\dagger$} & 230K \\
PLONK & Universal & No\textsuperscript{$\dagger$} & 350K \\
STARK & Transparent & Yes & 500K \\
\bottomrule
\end{tabular}
\caption{Production proof systems. $\dagger$: Wrapped in PQZ for quantum resistance.}
\end{table}

\subsubsection{Research Lane}

Experimental proof systems register under versioned capability IDs ($\geq 100$). Each carries an explicit risk label (\texttt{EXPERIMENTAL}, \texttt{AUDITED}, \texttt{DEPRECATED}) and runs in sandboxed verifier contracts that cannot affect production lane state.

\section{The PQZ Algorithm}

Post-Quantum Zero-Knowledge (PQZ) combines three primitives into a unified framework:

\subsection{Poseidon2 Commitments}

We replace Pedersen commitments (which rely on discrete log hardness) with Poseidon2 \cite{grassi2023poseidon2}, a ZK-friendly algebraic hash function secure under collision resistance assumptions that do not require elliptic curve hardness.

\begin{definition}[Poseidon2 Commitment]
For message $m \in \mathbb{F}_p^n$ and randomness $r \in \mathbb{F}_p$:
$$\text{Commit}(m, r) = \text{Poseidon2}(m \| r)$$
\end{definition}

\textbf{Performance.} Benchmarks on Zoo Network validators (NVIDIA H100):

\begin{table}[h]
\centering
\begin{tabular}{lrrr}
\toprule
\textbf{Operation} & \textbf{Poseidon2} & \textbf{Pedersen} & \textbf{Speedup} \\
\midrule
Commitment & 29 ns & 92 $\mu$s & 2000$\times$ \\
Note commit & 50 ns & 99 $\mu$s & 1980$\times$ \\
Merkle path (depth 32) & 928 ns & 2.9 ms & 3125$\times$ \\
Gas cost & 800 & 6,000 & 7.5$\times$ \\
\bottomrule
\end{tabular}
\caption{Poseidon2 vs Pedersen on GPU-accelerated validators.}
\end{table}

\subsection{STARK Proofs over Goldilocks}

For the core proving system, PQZ uses STARKs over the Goldilocks field ($p = 2^{64} - 2^{32} + 1$), which provides:

\begin{itemize}
\item \textbf{Transparency}: No trusted setup. Verification relies only on hash function collision resistance.
\item \textbf{Post-quantum security}: Security reduces to hash collision resistance (Poseidon2), not discrete log or factoring.
\item \textbf{Efficient FRI}: The Goldilocks field admits $2^{32}$-th roots of unity, enabling efficient Fast Reed-Solomon Interactive Oracle Proofs.
\end{itemize}

\subsection{ML-DSA Receipt Signing}

After STARK verification, the Z-Chain validator set signs the resulting receipt using ML-DSA (FIPS 204) threshold signatures coordinated through K-Chain:

\begin{enumerate}
\item Prover submits STARK proof to Z-Chain.
\item Z-Chain validators independently verify the proof.
\item Upon $\geq 67\%$ agreement, validators produce a threshold ML-DSA signature on the receipt.
\item The signed receipt is stored on-chain and indexed by receipt hash.
\end{enumerate}

\subsection{Groth16 Export for External EVMs}

For interoperability with Ethereum and other EVM chains that only support BN254 pairings, PQZ provides a Groth16 export:

\begin{enumerate}
\item Take a verified STARK receipt from Z-Chain.
\item Generate a Groth16 proof that ``I know a STARK proof that was verified and signed by the Zoo validator set.''
\item Export the Groth16 proof to the target EVM chain.
\end{enumerate}

This wraps the post-quantum STARK in a classical Groth16 for compatibility, while the original STARK receipt remains the source of truth on Z-Chain.

\section{Universal Receipt Format}

The core interoperability object:

\begin{definition}[Universal Receipt]
A receipt $R$ consists of:
\begin{align}
R = (&\texttt{programId},\ \texttt{claimHash},\ \texttt{receiptHash}, \notag \\
     &\texttt{proofSystemId},\ \texttt{version},\ \texttt{verifiedAt}, \notag \\
     &\texttt{parentReceipt},\ \texttt{aggregationRoot}) \notag
\end{align}
where:
\begin{itemize}
\item \texttt{programId}: Hash of the verified program/circuit.
\item \texttt{claimHash}: Hash of public inputs.
\item \texttt{receiptHash}: Self-referential hash ($H(R \setminus \texttt{receiptHash})$).
\item \texttt{proofSystemId}: $1$=STARK, $2$=Groth16, $3$=PLONK, $4$=Nova, $\geq 100$=research.
\item \texttt{parentReceipt}: For recursive composition.
\item \texttt{aggregationRoot}: Merkle root for batch receipts.
\end{itemize}
\end{definition}

\section{ZK Precompile ISA}

Z-Chain exposes cryptographic primitives as EVM precompiles at address range \texttt{0x0500--0x051F}:

\begin{table}[h]
\centering
\begin{tabular}{llr}
\toprule
\textbf{Address} & \textbf{Operation} & \textbf{Gas} \\
\midrule
\texttt{0x0501} & Poseidon2 hash & 800 \\
\texttt{0x0502} & Pedersen commit (legacy) & 6,000 \\
\texttt{0x0510} & STARK field arithmetic (Goldilocks) & 8--15 \\
\texttt{0x0511} & STARK extension field (Fp2) & 25 \\
\texttt{0x0512} & STARK Poseidon2 over Goldilocks & 800 \\
\texttt{0x0513} & STARK FRI fold & 2,000 \\
\texttt{0x0514} & STARK Merkle verification & 400/level \\
\texttt{0x051F} & STARK full proof verify & 100K--500K \\
\bottomrule
\end{tabular}
\caption{ZK Cryptographic ISA precompiles.}
\end{table}

\section{FHE-ZK Hybrid Protocols}

For private computation, Z-Chain integrates with T-Chain's threshold FHE:

\begin{enumerate}
\item \textbf{Encrypted input}: User encrypts data under Zoo's public FHE key.
\item \textbf{Homomorphic computation}: T-Chain validators evaluate the circuit on encrypted data.
\item \textbf{ZK proof of correctness}: A STARK proof attests that the computation was performed correctly on the ciphertext.
\item \textbf{Threshold decryption}: T-Chain validators (67-of-100 quorum) decrypt only the result, not the input.
\item \textbf{Receipt}: Z-Chain issues a receipt linking the FHE output to the ZK proof.
\end{enumerate}

This enables verifiable private inference: a model runs on encrypted patient data, the output is decrypted, and a ZK proof certifies correctness---all without exposing the data.

\section{Formal Verification}

All critical Z-Chain components are verified in Lean 4:

\begin{theorem}[STARK Soundness]
For any polynomial commitment scheme $\Pi$ over the Goldilocks field with security parameter $\lambda$, the probability of a computationally bounded adversary producing a valid STARK proof for a false statement is at most $2^{-\lambda}$.
\end{theorem}

\begin{theorem}[Receipt Integrity]
For any receipt $R$ with valid ML-DSA threshold signature from $\geq 67\%$ of validators, the probability that $R.\texttt{claimHash}$ does not correspond to a correctly verified proof is at most $2^{-128}$.
\end{theorem}

\begin{theorem}[PQZ Composition]
If the inner STARK proof is sound and the outer Groth16 wrapper is sound, then the composed PQZ proof is sound. The zero-knowledge property is preserved through both layers.
\end{theorem}

The full proofs are available at \texttt{zoo-formal/lean/ZK/} in the Zoo formal verification repository.

\section{Empirical Evaluation}

\subsection{Throughput}

On a cluster of 5 Zoo validators (each with NVIDIA H100 80GB):

\begin{table}[h]
\centering
\begin{tabular}{lrrr}
\toprule
\textbf{Operation} & \textbf{Throughput} & \textbf{Latency} & \textbf{Gas} \\
\midrule
STARK verify (simple) & 12,000/s & 1.1s & 100K \\
STARK verify (complex) & 3,200/s & 1.8s & 500K \\
Groth16 verify & 28,000/s & 0.4s & 230K \\
PLONK verify & 18,000/s & 0.6s & 350K \\
Receipt generation & 8,500/s & 1.4s & 50K \\
PQZ export (STARK$\to$Groth16) & 450/s & 8.2s & 800K \\
\bottomrule
\end{tabular}
\caption{Z-Chain verification throughput on GPU validators.}
\end{table}

\subsection{Comparison}

\begin{table}[h]
\centering
\begin{tabular}{lcccc}
\toprule
& \textbf{Z-Chain} & \textbf{=nil;} & \textbf{Polygon} & \textbf{StarkNet} \\
\midrule
PQ-safe & Yes & No & No & Partial \\
Multi-system & 3+ & 1 & 1 & 1 \\
Receipts & Native & No & No & No \\
FHE hybrid & Yes & No & No & No \\
GPU accel. & Yes & No & No & No \\
\bottomrule
\end{tabular}
\caption{Z-Chain vs existing ZK infrastructure.}
\end{table}

\section{Conclusion}

Z-Chain provides the first post-quantum zero-knowledge infrastructure integrated natively into a blockchain consensus layer. The PQZ algorithm unifies STARK transparency with ML-DSA quantum resistance and Groth16 EVM compatibility. Universal receipts enable composable proofs across chains. Formal verification in Lean 4 ensures soundness, completeness, and zero-knowledge. GPU acceleration achieves practical throughput (12,000 verifications/sec) for production deployment.

\begin{thebibliography}{20}

\bibitem{grassi2023poseidon2}
Grassi, L., Khovratovich, D., Rechberger, C., et al., ``Poseidon2: A New Hash Function for Zero-Knowledge Proof Systems,'' IACR ePrint 2023/323, 2023.

\bibitem{fips203}
NIST, ``Module-Lattice-Based Key-Encapsulation Mechanism Standard (ML-KEM),'' FIPS 203, 2024.

\bibitem{fips204}
NIST, ``Module-Lattice-Based Digital Signature Standard (ML-DSA),'' FIPS 204, 2024.

\bibitem{fips205}
NIST, ``Stateless Hash-Based Digital Signature Standard (SLH-DSA),'' FIPS 205, 2024.

\bibitem{stark}
Ben-Sasson, E., Bentov, I., Horesh, Y., Riabzev, M., ``Scalable, Transparent, and Post-Quantum Secure Computational Integrity,'' IACR ePrint 2018/046, 2018.

\bibitem{groth16}
Groth, J., ``On the Size of Pairing-Based Non-Interactive Arguments,'' EUROCRYPT 2016.

\bibitem{plonk}
Gabizon, A., Williamson, Z., Ciobotaru, O., ``PLONK: Permutations over Lagrange-Bases for Oecumenical Noninteractive Arguments of Knowledge,'' IACR ePrint 2019/953, 2019.

\bibitem{zoo-consensus}
Worring, A., Kelling, Z., ``Quasar Hybrid Consensus on Zoo Network,'' Zoo Labs Foundation, 2024.

\bibitem{zoo-pq}
Worring, A., Kelling, Z., ``Post-Quantum Cryptography Suite for Zoo Network,'' Zoo Labs Foundation, 2024.

\bibitem{zoo-fhe}
Worring, A., Kelling, Z., ``Threshold Fully Homomorphic Encryption on Zoo Network,'' Zoo Labs Foundation, 2025.

\end{thebibliography}

\end{document}
